% English translation of chapter7.tex lines 2036--2140.
% Source: Methods of Algebra, Volume 2, commit
% 9a5803ff77dd3257484cb177f851a73770a59dd3 (CC BY 4.0).
% stable-unit-id: o014.aljabr2.chapter7.identifying-module-categories
% Coverage: all of Section 7.8, ``Recognizing Module Categories''; ends
% immediately before Section 7.9 at line 2141.

% segment-id: o014.li2.chapter7.identifying-module-categories.h001
\section{Recognizing Module Categories}\label{sec:module-cat}
\hypertarget{o014.aljabr2.chapter7.identifying-module-categories}{ }
% segment-id: o014.li2.chapter7.identifying-module-categories.p001
Throughout this section, continue to fix a commutative ring $\Bbbk$. All
abelian categories and functors are understood to be $\Bbbk$-linear unless
otherwise stated.

% segment-id: o014.li2.chapter7.identifying-module-categories.p002
Categories of the form $R\dcate{Mod}$ or $\cated{Mod}R$, where $R$ is a
$\Bbbk$-algebra, are collectively called module categories. In
\S\ref{sec:Morita}, we specifically studied module categories and functors
between them. It is also natural to ask which categories are equivalent to
module categories and how such equivalences can be described concretely.
This section presents several concise, practical results about these
questions.

% segment-id: o014.li2.chapter7.identifying-module-categories.p003
Let $\mathcal{A}$ be an abelian category. Our approach is to consider a
functor of the form $\Hom_{\mathcal{A}}(s,\mathord\cdot)$, which naturally
lifts to a functor $\mathcal{A}\to\cated{Mod}R$ with
$R:=\End_{\mathcal{A}}(s)$. The question is when
$\Hom_{\mathcal{A}}(s,\mathord\cdot)$ is an equivalence.

\index{compact object}
% segment-id: o014.li2.chapter7.identifying-module-categories.p004
Recall Definition~\ref{def:cpt-objects} of a compact object. Suppose
$\mathcal{A}$ has all small filtered $\varinjlim$ and
$X\in\Obj(\mathcal{A})$. The object $X$ is called \emph{compact} if, for
every small filtered category $I$ and functor $\alpha:I\to\mathcal{A}$, the
following canonical morphism is an isomorphism:
% segment-id: o014.li2.chapter7.identifying-module-categories.q001
\[ \varinjlim_i \Hom\left( X, \alpha(i) \right) \to \Hom\left( X, \varinjlim \alpha \right). \]

% segment-id: o014.li2.chapter7.identifying-module-categories.le001
\begin{lemma}\label{prop:proj-compact}
	Let $X$ be a projective object in a cocomplete abelian category
	$\mathcal{A}$. Then $X$ is compact if and only if
	$\Hom_{\mathcal{A}}(X,\mathord\cdot)$ preserves all small direct sums;
	that is, for every small set $I$ and family of objects $(Y_i)_{i\in I}$,
	$\bigoplus_{i\in I}\Hom_{\mathcal{A}}(X,Y_i)\rightiso
	\Hom_{\mathcal{A}}(X,\bigoplus_{i\in I}Y_i)$.
	When this condition holds, $\Hom_{\mathcal{A}}(X,\mathord\cdot)$ in fact
	preserves all small $\varinjlim$.
\end{lemma}
% segment-id: o014.li2.chapter7.identifying-module-categories.pf001
\begin{proof}
	Suppose $X$ is compact. A direct sum indexed by a small set $I$ is the
	filtered $\varinjlim$ of the finite direct sums (take all finite subsets
	of $I$, ordered by inclusion). Since
	$\Hom_{\mathcal{A}}(X,\mathord\cdot)$ is known to preserve finite direct
	sums, it also preserves small direct sums.

	Conversely, projectivity is equivalent to
	$\Hom_{\mathcal{A}}(X,\mathord\cdot)$ preserving cokernels, while a
	general small $\varinjlim$ can be constructed from small direct sums and
	cokernels. Hence $\Hom_{\mathcal{A}}(X,\mathord\cdot)$ preserves all
	small $\varinjlim$.
\end{proof}

% segment-id: o014.li2.chapter7.identifying-module-categories.eg001
\begin{example}
	Let $R$ be a $\Bbbk$-algebra and $M$ a projective right $R$-module. Then
	$M$ is compact if and only if $M$ is finitely generated.
	\begin{itemize}
		\item Necessity: there are a small set $I$ and homomorphisms
		$i:M\to R^{\oplus I}$ and $p:R^{\oplus I}\to M$ with
		$pi=\identity_M$. Compactness implies that the image of $i$ is
		contained in $R^{\oplus I_0}$ for some finite subset $I_0\subset I$.
		Take $p':=p|_{R^{\oplus I_0}}$; then $p'i=\identity_M$ still holds.
		Thus $M$ can be realized as a direct summand of a finite-rank free
		module.
		\item Sufficiency: apply the criterion of
		Lemma~\ref{prop:proj-compact}. Every homomorphism
		$M\to\bigoplus_{i\in I}Y_i$ must land in
		$\bigoplus_{i\in I_0}Y_i$ for some finite subset $I_0\subset I$ (it
		is enough to check the images of the generators). Hence $M$ is compact.
	\end{itemize}
\end{example}

% segment-id: o014.li2.chapter7.identifying-module-categories.pr001
\begin{proposition}\label{prop:proj-compact-gen}
	Let $\mathcal{A}$ be a cocomplete abelian category. Then
	$s\in\Obj(\mathcal{A})$ is a compact projective generator if and only if
	$\Hom_{\mathcal{A}}(s,\mathord\cdot):
	\mathcal{A}\to\Bbbk\dcate{Mod}$ is a faithful exact functor preserving all
	small $\varinjlim$.
\end{proposition}
% segment-id: o014.li2.chapter7.identifying-module-categories.pf002
\begin{proof}
	By Proposition~\ref{prop:injective-cogenerator}, being a projective
	generator is equivalent to $\Hom_{\mathcal{A}}(s,\mathord\cdot)$ being
	faithful and exact. By Lemma~\ref{prop:proj-compact}, under the
	projectivity assumption, compactness is equivalent to
	$\Hom_{\mathcal{A}}(s,\mathord\cdot)$ preserving all small
	$\varinjlim$.
\end{proof}

% segment-id: o014.li2.chapter7.identifying-module-categories.th001
\begin{theorem}[P.\ Gabriel]\label{prop:identify-ModR}
	Let $s\in\Obj(\mathcal{A})$ and set
	$R:=\End_{\mathcal{A}}(s)$. Then
	% segment-id: o014.li2.chapter7.identifying-module-categories.q002
	\[ \Hom_{\mathcal{A}}(s, \cdot): \mathcal{A} \to \cated{Mod}R \]
	is an equivalence if and only if $s$ is a compact projective generator of
	$\mathcal{A}$.
\end{theorem}
% segment-id: o014.li2.chapter7.identifying-module-categories.pf003
\begin{proof}
	For brevity, write
	$G:=\Hom_{\mathcal{A}}(s,\mathord\cdot):\mathcal{A}\to\cated{Mod}R$.
	First prove necessity. Since $G$ is an equivalence and $G(s)=R$, it is
	enough to check that $R$ is a compact projective generator of
	$\cated{Mod}R$, which is immediate.

	We now prove sufficiency. Suppose $s$ is a compact projective generator.
	Proposition~\ref{prop:proj-compact-gen} implies that $G$ is a faithful
	exact functor preserving all small $\varinjlim$, because its composite with
	the forgetful functor $\cated{Mod}R\to\Bbbk\dcate{Mod}$ has these
	properties. Since faithfulness is already known, it is enough to prove
	that:
	\begin{compactitem}
		\item for all $X,Y\in\Obj(\mathcal{A})$, the corresponding map
		$\Hom_{\mathcal{A}}(X,Y)\to\Hom_R(GX,GY)$ is surjective;
		\item $G$ is essentially surjective.
	\end{compactitem}

	First, for every $X\in\Obj(\mathcal{A})$, define a morphism
	$\phi_X:s^{\oplus\Hom_{\mathcal{A}}(s,X)}\to X$ whose restriction to the
	direct summand corresponding to $f:s\to X$ is $f$. We show that $\phi_X$
	is surjective. Otherwise, suppose that $\Coker(\phi_X)\neq0$. The generator
	property gives a nonzero morphism $s\to\Coker(\phi_X)$
	(Proposition~\ref{prop:injective-cogenerator}), while projectivity makes
	this morphism factor as $s\xrightarrow{g}X\to\Coker(\phi_X)$. It follows
	that the composite
	$ s^{\oplus\Hom_{\mathcal{A}}(s,X)}\xrightarrow{\phi_X}X
	\to\Coker(\phi_X)$ is nonzero on the direct summand corresponding to $g$,
	a contradiction.

	Apply the same construction to the kernel $\Ker(\phi_X)$. Thus every $X$
	fits into an exact sequence
	% segment-id: o014.li2.chapter7.identifying-module-categories.q003
	\[ s^{\oplus J} \to s^{\oplus I} \to X \to 0, \quad I = I_X, \; J = J_X: \text{small sets}. \]
	Since $G$ preserves all small $\varinjlim$, this gives an exact sequence of
	right $R$-modules
	% segment-id: o014.li2.chapter7.identifying-module-categories.q004
	\[ R^{\oplus J} \to R^{\oplus I} \to GX \to 0. \]
	The first morphism in each sequence may be regarded as an $I\times J$
	matrix over $R$ acting by left multiplication; the two are in fact the
	same matrix.

	Now consider a homomorphism of right $R$-modules $f:GX\to GY$. Choose
	exact sequences of the preceding form for $X$ and $Y$. It is easy to see
	that $f$ lifts to a commutative diagram with exact rows
	% segment-id: o014.li2.chapter7.identifying-module-categories.q005
	\[\begin{tikzcd}
		R^{\oplus J_X} \arrow[r] \arrow[d] & R^{\oplus I_X} \arrow[d] \arrow[r] & GX \arrow[d, "f"] \arrow[r] & 0 \\
		R^{\oplus J_Y} \arrow[r] & R^{\oplus I_Y} \arrow[r] & GY \arrow[r] & 0.
	\end{tikzcd}\]
	The four arrows in the left square, however, can all be expressed as
	matrices over $R$. Those four matrices give the solid part of a
	commutative diagram with exact rows
	% segment-id: o014.li2.chapter7.identifying-module-categories.q006
	\[\begin{tikzcd}
		s^{\oplus J_X} \arrow[r] \arrow[d] & s^{\oplus I_X} \arrow[d] \arrow[r] & X \arrow[dashed, d] \arrow[r] & 0 \\
		s^{\oplus J_Y} \arrow[r] & s^{\oplus I_Y} \arrow[r] & Y \arrow[r] & 0,
	\end{tikzcd}\]
	and the dashed part is uniquely determined by functoriality of the
	cokernel. Since $G$ is exact, necessarily $G(X\dashrightarrow Y)=f$.

	A similar argument shows that $G$ is essentially surjective. For a given
	right $R$-module $M$, there is an exact sequence
	% segment-id: o014.li2.chapter7.identifying-module-categories.q007
	\[ R^{\oplus L} \to R^{\oplus K} \to M \to 0, \]
	where the first morphism is regarded as a $K\times L$ matrix over $R$.
	The same matrix gives a cokernel exact sequence in $\mathcal{A}$,
	% segment-id: o014.li2.chapter7.identifying-module-categories.q008
	\[ s^{\oplus L} \to s^{\oplus K} \to X \to 0. \]
	Since $G$ preserves all small $\varinjlim$, necessarily $GX\simeq M$.
	This completes the proof.
\end{proof}

% segment-id: o014.li2.chapter7.identifying-module-categories.p005
Another related problem is to recognize categories of finitely generated
modules. Write $\catesubd{Mod}{\mathrm{fg}}R$ for the category of finitely
generated right $R$-modules. If $R$ is a right Noetherian ring,
$\catesubd{Mod}{\mathrm{fg}}R$ remains an abelian category; see
\cite[Proposition 6.10.3, Lemma 6.10.4]{Li1}.
\index[sym1]{Modfg@$\cate{Mod}_{\mathrm{fg}}$}

% segment-id: o014.li2.chapter7.identifying-module-categories.th002
\begin{theorem}\label{prop:identify-ModR-fg}
	Let $\mathcal{A}$ be an abelian category and $s\in\Obj(\mathcal{A})$.
	Suppose that:
	\begin{compactitem}
		\item every $X\in\Obj(\mathcal{A})$ is a Noetherian object
		(Definition~\ref{def:finite-length-object});
		\item $R:=\End_{\mathcal{A}}(s)$ is a right Noetherian ring;
		\item $s$ is a projective generator of $\mathcal{A}$.
	\end{compactitem}
	Then the functor $\Hom_{\mathcal{A}}(s,\mathord\cdot):
	\mathcal{A}\to\cated{Mod}R$ factors as
	% segment-id: o014.li2.chapter7.identifying-module-categories.q009
	\[ \mathcal{A} \xrightarrow{\text{equivalence}} \catesubd{Mod}{\mathrm{fg}}R \subset \cated{Mod}R. \]
\end{theorem}
% segment-id: o014.li2.chapter7.identifying-module-categories.pf004
\begin{proof}
	First show that every $X\in\Obj(\mathcal{A})$ can be placed in an exact
	sequence
	% segment-id: o014.li2.chapter7.identifying-module-categories.q010
	\begin{equation}\label{eqn:identify-ModR-fg-aux}
		s^{\oplus m} \to s^{\oplus n} \to X \to 0, \quad n, m \in \Z_{\geq 0}.
	\end{equation}

	Indeed, since $s$ is a generator, for every $X\neq0$ there is a nonzero
	morphism $\phi_1:s\to X$. If $\phi_1$ is not surjective, there is a
	nonzero morphism $s\to\Coker(\phi_1)$; the projectivity of $s$ ensures
	that this morphism factors as
	$s\xrightarrow{\psi_1}X\to\Coker(\phi_1)$. We thereby obtain
	$\phi_2:=(\phi_1,\psi_1):s^{\oplus2}\to X$ with
	$\Image(\phi_2)\supsetneq\Image(\phi_1)$. If $\phi_2$ is not yet
	surjective, repeat the same procedure. Since $X$ is a Noetherian object,
	this process must terminate after finitely many steps and yield a
	surjective morphism $\phi_X:s^{\oplus n}\twoheadrightarrow X$. Apply the
	same procedure to $\Ker(\phi_X)$ to obtain
	\eqref{eqn:identify-ModR-fg-aux}.

	Recall that $s$ is a projective generator if and only if
	$G:=\Hom_{\mathcal{A}}(s,\mathord\cdot)$ is faithful and exact. Applying
	$G$ to \eqref{eqn:identify-ModR-fg-aux} gives an exact sequence
	$R^{\oplus m}\to R^{\oplus n}\to GX\to0$, showing that the image of $G$
	lies in $\catesubd{Mod}{\mathrm{fg}}R$.

	On the other hand, since $R$ is right Noetherian, every
	$M\in\Obj(\catesubd{Mod}{\mathrm{fg}}R)$ can also be placed in an exact
	sequence
	% segment-id: o014.li2.chapter7.identifying-module-categories.q011
	\[ R^{\oplus q} \to R^{\oplus p} \to M \to 0, \quad p, q \in \Z_{\geq 0}. \]

	The rest of the proof follows the same idea as the proof of sufficiency in
	Theorem~\ref{prop:identify-ModR}. In that argument, the compactness of $s$
	was used only to handle infinite direct sums; here only finite direct sums
	occur, and every additive functor preserves them. The remainder of the
	argument is identical and is not repeated.
\end{proof}
